% !TEX output_directory = ../publica
\documentclass[twoside, 10pt, a4paper]{article}
\usepackage[utf8]{inputenc}
\usepackage[T1]{fontenc}
\usepackage[portuguese]{babel}

% --- FONTE PADRÃO (Estilo Didático Encorpado) ---
\usepackage{helvet} 
\renewcommand{\familydefault}{\sfdefault}

% --- GEOMETRIA E ESPAÇAMENTO ---
\usepackage[top=1.6cm, bottom=1.5cm, left=0.9cm, right=0.9cm, headheight=24pt, headsep=0.4cm]{geometry}
\usepackage{microtype} 
\usepackage{setspace}
\setstretch{1.0} % Espaçamento simples
\usepackage{parskip}
\setlength{\parskip}{2pt}
\setlength{\parindent}{0pt}

% ==========================================
%   PAINEL DE CONTROLE DE ESPAÇAMENTO E ESCALA
% ==========================================
\newlength{\espacoPosTitulo}
\setlength{\espacoPosTitulo}{0.3cm}  % Espaço entre o título do capítulo e a 1ª questão

\newlength{\espacoPreQuestao}
\setlength{\espacoPreQuestao}{0.1cm} % Espaço entre a linha fina divisória e a questão

\newlength{\espacoQuestao}
\setlength{\espacoQuestao}{0.2cm}    % Espaço de "respiro" no FINAL de cada questão

\newlength{\espacoAlt}
\setlength{\espacoAlt}{0.2cm}        % Espaço para quebras manuais de alternativas

% Controle de Listas Automáticas (Ambientes enumerate / itemize)
\usepackage{enumitem}
\setlist{
    itemsep=0.4cm,    
    topsep=0.1cm,     
    parsep=0pt,       
    partopsep=0pt     
}

% Controle de Escala de Imagens
\newcommand{\larguraTikz}{0.85\linewidth} 
% ==========================================

% --- MATEMÁTICA ---
\usepackage{amsmath, amsfonts, amssymb}
\usepackage{nccmath}
\usepackage{mathastext}

% --- CORES CUSTOMIZADAS PREMIUM ---
\usepackage{xcolor}
\definecolor{indigo}{HTML}{4F46E5}
\definecolor{CorTitUm}{HTML}{FF6F61}
\definecolor{CorTitDois}{HTML}{FF6F61}
\definecolor{CorTitTres}{HTML}{658894}
\definecolor{CorLinha}{HTML}{B0B0B0}   
\definecolor{metal}{HTML}{7F8C8D}
\definecolor{vidro}{HTML}{3498DB}
\definecolor{ouro}{HTML}{F1C40F}
\definecolor{concreto}{HTML}{95A5A6}
\definecolor{madeira}{HTML}{D35400}
\definecolor{CinzaEscuro}{HTML}{4A4A4A} 
\definecolor{CinzaMedio}{HTML}{848484} 
\definecolor{Cinzablack}{HTML}{353535} 

% --- MINI-CABEÇALHO E RODAPÉ AUTORAL (TWOSIDE ESPELHADO) ---
\usepackage{fancyhdr}
\pagestyle{fancy}
\fancyhf{} 
\renewcommand{\headrulewidth}{0.3pt} 
\renewcommand{\footrulewidth}{0.3pt}
\fancyhead[LE]{\small\sffamily\bfseries\color{gray!75} ESTRUTURA DA TERRA E ATMOSFERA}
\fancyhead[RO]{\small\sffamily\color{gray!75} \texttt{profraphaelpaulino.com.br}}
\fancyfoot[LE]{\small \textbf{\thepage}}
\fancyfoot[RO]{\small \textbf{\thepage}}

% --- CAIXAS DE DESTAQUE (tcolorbox) ---
\usepackage[most]{tcolorbox}

% 1. Box de Resposta do Aluno
\newtcolorbox{boxresponda}{
    colback=white,        
    colframe=gray!45,       
    boxrule=0pt,          
    leftrule=1.7pt,        
    sharp corners,        
    enlarge left by=0.4cm,  
    width=\linewidth-0.4cm, 
    left=8pt, right=0pt,  
    top=2pt, bottom=2pt,
    fontupper=\small 
}

% 2. Box de Dica (Estilo Ciano Claríssimo)
\definecolor{CorDicaBorda}{HTML}{17c1be} 
\definecolor{CorDicaFundo}{HTML}{f3fbfb} 
\newtcolorbox{boxdica}{
    colback=CorDicaFundo,        
    colframe=CorDicaBorda,      
    boxrule=0pt,          
    leftrule=2.5pt,         
    sharp corners,        
    width=\linewidth, 
    left=8pt, right=8pt,  
    top=6pt, bottom=6pt,
    halign=center, 
    fontupper=\small      
}

% --- TÍTULOS (ESTILO EDITORIAL) ---
\usepackage{titlesec}
\titleformat{\section}{\color{black}\large\bfseries}{\thesection}{0em}{\vspace{0.1cm}\titlerule[0.8pt]}
\titlespacing*{\section}{0pt}{10pt}{6pt} 

\titleformat{\subsubsection}[runin]{\color{black}\bfseries}{}{0em}{}
\titlespacing*{\subsubsection}{0pt}{0pt}{0.15cm}

% --- PACOTES DE ESTRUTURA E GRÁFICOS ---
\usepackage{multicol}
\setlength{\columnsep}{1cm}
\setlength{\columnseprule}{0.5pt}
\def\columnseprulecolor{\color{CorLinha}}

\usepackage{adjustbox}
\usepackage{tikz}
\usepackage{pgfplots}
\usepackage{circuitikz}
\usetikzlibrary{shadows, shapes.misc, positioning, calc}
\pgfplotsset{compat=1.18}

\begin{document}
\thispagestyle{empty}
\color{Cinzablack}
\vspace*{-1.8cm}

% --- CABEÇALHO PRINCIPAL AUTORAL (Apenas 1ª página) ---
\noindent
\begin{minipage}{\textwidth}
    \fontfamily{phv}\selectfont
    \noindent
    \begin{minipage}[t]{0.70\linewidth}
        {\Large \bfseries \MakeUppercase{Caderno de Atividades Suplementares}}\\[0.1cm]
        {\large \textmd{\textcolor{CinzaEscuro}{Unidade: Estrutura da Terra e Atmosfera}}}
    \end{minipage}%
    \begin{minipage}[t]{0.30\linewidth}
        \raggedleft
        \small \textbf{Prof. Raphael Paulino}\\
        \textsf{\small \textcolor{indigo}{\texttt{profraphaelpaulino.com.br}}}
    \end{minipage}
    
    \vspace{0.15cm}
    \noindent\rule{\linewidth}{1.2pt}
    \vspace{-0.1cm}
    \footnotesize\textsf{\textbf{ÁREA:} Física \hfill \textbf{NÍVEL:} EF II (6º Ano)}
    \vspace{0.1cm}
    \noindent\rule{\linewidth}{0.4pt}
\end{minipage}

\par\vspace{0.3cm} 
\raggedcolumns
\begin{multicols*}{2}
\vspace{0.1cm}

\subsection*{Capítulo 1: Camadas Internas da Terra}
\vspace{-0.2cm}\noindent\rule{\linewidth}{1.5pt} 
\par\vspace{\espacoPosTitulo}
\subsubsection*{01.} O diagrama ilustra um modelo do interior do nosso planeta, dividido didaticamente em quatro grandes regiões estruturais, baseadas em composição química e estado físico.

\begin{center}
% ==========================================
% CONTROLE INDIVIDUAL DESTA IMAGEM:
% 1. CORTAR BORDAS: Mude os zeros do trim (Esquerda, Baixo, Direita, Topo). Ex: trim=1cm 0cm 1cm 0cm
% 2. TAMANHO: Para encolher só esta imagem, troque \larguraTikz por um valor (Ex: 0.5\linewidth)
% ==========================================
\begin{adjustbox}{trim=0cm 0cm 0cm 0cm, clip, max width=\larguraTikz}
\begin{tikzpicture}
\definecolor{crosta}{HTML}{8D6E63}
\definecolor{manto}{HTML}{E67E22}
\definecolor{nucleoext}{HTML}{F1C40F}
\definecolor{nucleoint}{HTML}{E74C3C}
\definecolor{fundo}{HTML}{F8F9FA}

% --- APLICANDO DROP SHADOW ---
\fill[fundo, rounded corners=5pt, drop shadow={opacity=0.15, shadow xshift=2pt, shadow yshift=-2pt}] (-3.5,-3.5) rectangle (3.5,3.5);

\fill[crosta, drop shadow={opacity=0.2}] (0,0) circle (2.8);
\fill[manto] (0,0) circle (2.6);
\fill[nucleoext] (0,0) circle (1.5);
\fill[nucleoint] (0,0) circle (0.7);

\fill[fundo] (0,0) -- (3,0) arc (0:45:3) -- cycle;
\draw[thick, white] (0,0) -- (2.8,0);
\draw[thick, white] (0,0) -- (1.98,1.98);

\node[fill=fundo, inner sep=2pt, rounded corners=2pt, font=\sffamily\footnotesize, align=center] at (-1.3,2) {Região B};
\node[fill=fundo, inner sep=2pt, rounded corners=2pt, font=\sffamily\footnotesize, align=center] at (-0.8,0.8) {Região C};
\node[fill=fundo, inner sep=2pt, rounded corners=2pt, font=\sffamily\footnotesize, align=center] at (0,-0.3) {Região D};
\draw[stealth-] (2.7, -0.5) -- (3.2, -1) node[fill=fundo, right, inner sep=2pt, rounded corners=2pt, font=\sffamily\footnotesize, align=center] {Região A};
\end{tikzpicture}
\end{adjustbox}
\end{center}

\begin{boxresponda}
\textbf{Responda:} Com base na dinâmica terrestre, identifique qual região (A, B, C ou D) corresponde à camada onde ocorrem os movimentos convectivos do magma e explique como esse movimento interfere na crosta terrestre.
\end{boxresponda}
\par\vspace{\espacoQuestao}

\noindent\textcolor{CorLinha}{\rule{\linewidth}{0.5pt}} 
\par\vspace{\espacoPreQuestao}
\subsubsection*{02.} A compreensão das camadas da Terra envolve analisar não apenas a profundidade, mas o estado físico dos materiais. Julgue as afirmativas a seguir, classificando-as como Verdadeiras (V) ou Falsas (F).

\begin{enumerate}[label=(\ \ ), itemsep=\espacoAlt]
\item O núcleo interno, apesar das altíssimas temperaturas, encontra-se no estado sólido devido à imensa pressão exercida pelas camadas superiores.
\item A crosta terrestre é uma camada contínua e uniforme, possuindo a mesma espessura tanto sob os oceanos quanto sob os continentes.
\item O núcleo externo é líquido e rico em ferro e níquel, sendo o principal responsável pela geração do campo magnético da Terra.
\item O manto é formado unicamente por água em estado de ebulição, o que gera as erupções vulcânicas.
\end{enumerate}
\par\vspace{\espacoQuestao}

\noindent\textcolor{CorLinha}{\rule{\linewidth}{0.5pt}} 
\par\vspace{\espacoPreQuestao}
\subsubsection*{03.} As rochas que compõem a litosfera estão sujeitas a diferentes condições dependendo da profundidade. Justifique por que o núcleo interno, mesmo apresentando temperaturas que podem ultrapassar os \textbf{5000 \textdegree C}, não derrete e permanece no estado sólido.
\par\vspace{\espacoQuestao}

\noindent\textcolor{CorLinha}{\rule{\linewidth}{0.5pt}} 
\par\vspace{\espacoPreQuestao}
\subsubsection*{04.} Os poços de perfuração mais profundos já escavados pela humanidade atingem pouco mais de \textbf{12 km} de profundidade, uma fração minúscula em relação ao raio do nosso planeta, que é de aproximadamente \textbf{6370 km}. Demonstre como os cientistas conseguem mapear e determinar o estado físico de camadas tão profundas sem nunca terem chegado fisicamente até lá.
\par\vspace{\espacoQuestao}

\noindent\textcolor{CorLinha}{\rule{\linewidth}{0.5pt}} 
\par\vspace{\espacoPreQuestao}
\subsubsection*{05.} (Estilo OBA) O sismograma é um registro gráfico gerado pela propagação de energia originada por abalos sísmicos. O esquema apresenta o comportamento das ondas S (secundárias), que possuem a propriedade de não se propagarem em meios líquidos.

\begin{center}
% ==========================================
% CONTROLE INDIVIDUAL DESTA IMAGEM:
% ==========================================
\begin{adjustbox}{trim=0cm 0cm 0cm 0cm, clip, max width=\larguraTikz}
\begin{tikzpicture}
\definecolor{terracota}{HTML}{D35400}
\definecolor{cinzamedio}{HTML}{7F8C8D}
\definecolor{fundo}{HTML}{FDFEFE}

\fill[fundo, rounded corners=5pt, drop shadow={opacity=0.15, shadow xshift=2pt, shadow yshift=-2pt}] (-3,-3) rectangle (3,3);

\fill[terracota!20] (0,0) circle (2.5);
\draw[thick, cinzamedio] (0,0) circle (2.5);
\fill[terracota!50] (0,0) circle (1.3);
\draw[dashed, cinzamedio] (0,0) circle (1.3);
\fill[terracota!80] (0,0) circle (0.6);

\fill[red] (0,2.5) circle (0.08);
\node[fill=fundo, inner sep=1pt, rounded corners=1pt, font=\sffamily\tiny, align=center] at (0, 2.7) {Epicentro};

\draw[thick, -stealth] (0,2.5) -- (-1.5, 1);
\draw[thick, -stealth] (0,2.5) -- (1.5, 1);
\draw[thick, -stealth] (0,2.5) -- (-2.2, 0);
\draw[thick, -stealth] (0,2.5) -- (2.2, 0);

\draw[thick, -stealth] (0,2.5) -- (-0.6, 1.15);
\node[red, font=\huge, align=center] at (-0.6, 1.15) {x};
\draw[thick, -stealth] (0,2.5) -- (0.6, 1.15);
\node[red, font=\huge, align=center] at (0.6, 1.15) {x};

\node[fill=fundo, inner sep=2pt, font=\sffamily\footnotesize, align=center] at (0,-1.8) {Zona de Sombra \\ (Sem registro de ondas S)};
\end{tikzpicture}
\end{adjustbox}
\end{center}

Analisando a trajetória geométrica das frentes de onda ilustradas e sabendo das propriedades mecânicas das ondas S, conclua qual é o estado físico da camada imediatamente abaixo do manto terrestre.
\par\vspace{\espacoQuestao}

\noindent\textcolor{CorLinha}{\rule{\linewidth}{0.5pt}} 
\par\vspace{\espacoPreQuestao}
\subsubsection*{06.} (Estilo VUNESP) Em diversas obras de ficção científica, como no clássico "Viagem ao Centro da Terra" de Júlio Verne, personagens exploram cavernas colossais repletas de cristais e oceanos subterrâneos. Do ponto de vista científico moderno, uma jornada real ao centro do planeta por seres humanos seria estruturalmente impossível. 

Qual fator físico e geológico torna essa premissa ficcional impraticável para a engenharia atual?

\begin{enumerate}[label=\alph*), itemsep=\espacoAlt]
\item A ausência de minerais metálicos no núcleo terrestre, o que inviabiliza a comunicação de rádio.
\item A diminuição contínua da temperatura, que faria o maquinário congelar antes de atingir o manto.
\item O aumento extremo e simultâneo de temperatura e pressão, que derreteria e esmagaria qualquer sonda ou veículo.
\item A presença de bolsões de vácuo no manto superior, impedindo a locomoção de veículos escavadores.
\item A força da gravidade que atua de baixo para cima, expulsando materiais metálicos para a superfície.
\end{enumerate}
\par\vspace{\espacoQuestao}

\subsection*{Capítulo 2: Atmosfera}
\vspace{-0.2cm}\noindent\rule{\linewidth}{1.5pt} 
\par\vspace{\espacoPosTitulo}
\subsubsection*{07.} O gradiente vertical da atmosfera exibe diferentes estratificações químicas e térmicas. A figura representa a elevação quilométrica a partir do nível do mar.

\begin{center}
% ==========================================
% CONTROLE INDIVIDUAL DESTA IMAGEM:
% ==========================================
\begin{adjustbox}{trim=0cm 0cm 0cm 0cm, clip, max width=\larguraTikz}
\begin{tikzpicture}
\definecolor{tropo}{HTML}{85C1E9}
\definecolor{estrato}{HTML}{5DADE2}
\definecolor{meso}{HTML}{2E86C1}
\definecolor{termo}{HTML}{1B4F72}
\definecolor{fundo}{HTML}{F4F6F6}

\fill[fundo, rounded corners=5pt, drop shadow={opacity=0.15, shadow xshift=2pt, shadow yshift=-2pt}] (-0.5,-0.5) rectangle (6,5);

\fill[tropo, opacity=0.8] (0,0) rectangle (4,1);
\fill[estrato, opacity=0.8] (0,1) rectangle (4,2.2);
\fill[meso, opacity=0.8] (0,2.2) rectangle (4,3.4);
\fill[termo, opacity=0.8] (0,3.4) rectangle (4,4.5);

\node[font=\sffamily\footnotesize, white, align=center] at (2,0.5) {Camada 1};
\node[font=\sffamily\footnotesize, white, align=center] at (2,1.6) {Camada 2};
\node[font=\sffamily\footnotesize, white, align=center] at (2,2.8) {Camada 3};
\node[font=\sffamily\footnotesize, white, align=center] at (2,3.95) {Camada 4};

\draw[thick, -stealth] (4.2,0) -- (4.2,4.8);
\node[font=\sffamily\tiny, align=center, rotate=90] at (4.5,2.4) {Altitude (km)};

\node[font=\sffamily\tiny, align=center] at (4.8,1) {\textbf{15 km}};
\node[font=\sffamily\tiny, align=center] at (4.8,2.2) {\textbf{50 km}};
\node[font=\sffamily\tiny, align=center] at (4.8,3.4) {\textbf{85 km}};

\fill[white, opacity=0.7] (0.5,0.7) circle (0.15);
\fill[white, opacity=0.7] (0.7,0.8) circle (0.2);
\fill[white, opacity=0.7] (0.9,0.7) circle (0.15);

\draw[white, dashed, thick] (0,1.9) -- (4,1.9);
\node[fill=estrato, inner sep=1pt, font=\sffamily\tiny, white, align=center] at (2,1.9) {$O_3$};
\end{tikzpicture}
\end{adjustbox}
\end{center}

Com base na disposição altimétrica, especifique em qual camada (1, 2, 3 ou 4) ocorrem as precipitações meteorológicas (chuvas e furacões) e em qual camada localiza-se o escudo gasoso que nos protege da radiação ultravioleta.
\par\vspace{\espacoQuestao}

\noindent\textcolor{CorLinha}{\rule{\linewidth}{0.5pt}} 
\par\vspace{\espacoPreQuestao}
\subsubsection*{08.} O Efeito Estufa é frequentemente noticiado de forma negativa devido ao aquecimento global. Contudo, em sua essência natural, ele é um mecanismo vital. Demonstre a importância da atmosfera e do ciclo térmico para a manutenção da temperatura média global e a preservação da água no estado líquido.
\par\vspace{\espacoQuestao}

\noindent\textcolor{CorLinha}{\rule{\linewidth}{0.5pt}} 
\par\vspace{\espacoPreQuestao}
\subsubsection*{09.} As zonas superiores da atmosfera apresentam características únicas ligadas à densidade gasosa e à radiação solar. Julgue as afirmativas como Verdadeiras (V) ou Falsas (F).

\begin{enumerate}[label=(\ \ ), itemsep=\espacoAlt]
\item A Mesosfera é conhecida por registrar as temperaturas mais baixas da atmosfera e é onde a maioria dos meteoroides se desintegram por atrito.
\item Na Termosfera, o ar é tão denso que dificulta o voo de aeronaves comerciais e jatos.
\item A Exosfera representa a última fronteira antes do vácuo espacial, onde os gases se dissipam e satélites artificiais orbitam o planeta.
\item As auroras boreais são fenômenos luminosos que ocorrem exclusivamente na Troposfera, junto às nuvens de tempestade.
\end{enumerate}
\par\vspace{\espacoQuestao}

\noindent\textcolor{CorLinha}{\rule{\linewidth}{0.5pt}} 
\par\vspace{\espacoPreQuestao}
\subsubsection*{10.} Organize as ideias: Associe diretamente cada um dos quatro fenômenos listados a seguir com a sua respectiva camada atmosférica de ocorrência predominante. 

\begin{enumerate}[label=\Roman*., itemsep=\espacoAlt]
\item Formação de Granizo e Ciclones.
\item Absorção de Raios Ultravioleta.
\item Atrito e fragmentação de rochas espaciais (Estrelas Cadentes).
\item Interação de ventos solares com partículas magnéticas (Auroras).
\end{enumerate}
\par\vspace{\espacoQuestao}

\noindent\textcolor{CorLinha}{\rule{\linewidth}{0.5pt}} 
\par\vspace{\espacoPreQuestao}
\subsubsection*{11.} A pressão exercida por uma coluna de gás sofre variações dramáticas dependendo do referencial topográfico. O gráfico correlaciona essa pressão com a elevação do terreno.

\begin{center}
% ==========================================
% CONTROLE INDIVIDUAL DESTA IMAGEM:
% ==========================================
\begin{adjustbox}{trim=0cm 0cm 0cm 1cm, clip, max width={0.75\linewidth}}
\begin{tikzpicture}
\definecolor{azulmar}{HTML}{2980B9}
\definecolor{fundo}{HTML}{F9EBEA}

\fill[fundo, rounded corners=5pt, drop shadow={opacity=0.15, shadow xshift=2pt, shadow yshift=-2pt}] (-1,-1) rectangle (6.5,5);

\begin{axis}[
    width=6cm, height=5cm,
    axis lines=middle,
    xlabel={Altitude (km)},
    ylabel={Pressão},
    xmin=0, xmax=10,
    ymin=0, ymax=1.2,
    xtick={0, 2, 4, 6, 8, 8.8},
    xticklabels={0, 2, 4, 6, 8, \textbf{Everest}},
    ytick=\empty,
    every axis x label/.style={at={(ticklabel* cs:1.05)}, anchor=west, font=\footnotesize, align=center},
    every axis y label/.style={at={(ticklabel* cs:1.05)}, anchor=south, font=\footnotesize, align=center},
    tick label style={font=\footnotesize},
    x tick label style={rotate=45, anchor=east}
]
\addplot [domain=0:9, samples=100, thick, azulmar, drop shadow={opacity=0.3}] {exp(-0.15*x)};
\addplot [dashed, gray] coordinates {(8.8,0) (8.8, 0.267)};
\addplot [mark=*, mark size=2pt, red] coordinates {(8.8, 0.267)};
\end{axis}
\end{tikzpicture}
\end{adjustbox}
\end{center}

\begin{boxresponda}
\textbf{Responda:} Baseado na interpretação gráfica do decaimento pressórico, justifique por que alpinistas que escalam montanhas acima de \textbf{8000 m} necessitam utilizar cilindros de oxigênio suplementar, mesmo sabendo que a proporção percentual de oxigênio no ar continua sendo de aproximadamente \textbf{21\%}.
\end{boxresponda}
\par\vspace{\espacoQuestao}

\noindent\textcolor{CorLinha}{\rule{\linewidth}{0.5pt}} 
\par\vspace{\espacoPreQuestao}
\subsubsection*{12.} (ENEM - Adaptada) A ação humana tem gerado impactos severos na dinâmica dos gases do nosso planeta. A emissão histórica de compostos como os CFCs (Clorofluorcarbonetos) desencadeou um processo de degradação em uma importante camada de proteção terrestre, elevando os índices de doenças de pele e cataratas na população global. 

O escudo gasoso diretamente afetado por esses compostos e a camada da atmosfera onde ele se encontra alojado são, respectivamente:

\begin{enumerate}[label=\alph*), itemsep=\espacoAlt]
\item Camada de Metano na Mesosfera.
\item Camada de Ozônio na Estratosfera.
\item Cinturão Magnético na Exosfera.
\item Camada de Nitrogênio na Troposfera.
\item Gás Carbônico na Termosfera.
\end{enumerate}
\par\vspace{\espacoQuestao}
\subsection*{Capítulo 3: Rochas e Minerais}
\vspace{-0.2cm}\noindent\rule{\linewidth}{1.5pt} 
\par\vspace{\espacoPosTitulo}
\subsubsection*{13.} O fluxograma ilustra o modelo geológico contínuo de transformação dos materiais da litosfera terrestre. As setas representam processos físicos e químicos que ocorrem ao longo de milhões de anos.

\begin{center}
% ==========================================
% CONTROLE INDIVIDUAL DESTA IMAGEM:
% ==========================================
\begin{adjustbox}{trim=0cm 1.1cm 0cm 1.4cm, clip, max width={0.75\linewidth}}
\begin{tikzpicture}
\definecolor{igneas}{HTML}{E74C3C}
\definecolor{sedimentares}{HTML}{F1C40F}
\definecolor{metamorficas}{HTML}{9B59B6}
\definecolor{magma}{HTML}{D35400}
\definecolor{fundo}{HTML}{F8F9FA}

\fill[fundo, rounded corners=5pt, drop shadow={opacity=0.15, shadow xshift=2pt, shadow yshift=-2pt}] (-4.5,-3.8) rectangle (4.5,4.5);

\node[fill=igneas, text=white, font=\sffamily\bfseries, rounded corners=3pt, minimum width=2.4cm, minimum height=1cm, drop shadow={opacity=0.2}] (R1) at (-2.5,2) {Rochas Ígneas};
\node[fill=sedimentares, text=black, font=\sffamily\bfseries, rounded corners=3pt, minimum width=2.8cm, minimum height=1cm, drop shadow={opacity=0.2}] (R2) at (2.5,2) {Rochas Sedimentares};
\node[fill=metamorficas, text=white, font=\sffamily\bfseries, rounded corners=3pt, minimum width=2.8cm, minimum height=1cm, drop shadow={opacity=0.2}] (R3) at (1.5,-1.5) {Rochas Metamórficas};
\node[fill=magma, text=white, font=\sffamily\bfseries, rounded corners=3pt, minimum width=1.8cm, minimum height=0.8cm, drop shadow={opacity=0.2}] (Magma) at (-2.5,-1.5) {Magma};

\draw[-stealth, thick, color=gray] (Magma) -- (R1) node[midway, left, fill=fundo, inner sep=2pt, font=\sffamily\tiny, align=center] {Resfriamento};
\draw[-stealth, thick, color=gray] (R1) -- (R2) node[midway, above, fill=fundo, inner sep=2pt, font=\sffamily\tiny, align=center] {Processo X};
\draw[-stealth, thick, color=gray] (R2) -- (R3) node[midway, right, fill=fundo, inner sep=2pt, font=\sffamily\tiny, align=center] {Pressão e \\ Temp.};
\draw[-stealth, thick, color=gray] (R3) -- (Magma) node[midway, below, fill=fundo, inner sep=2pt, font=\sffamily\tiny, align=center] {Fusão};
\draw[-stealth, thick, color=gray] (R1) -- (R3) node[midway, right, fill=fundo, inner sep=2pt, font=\sffamily\tiny, align=center] {Calor e Pressão};
\end{tikzpicture}
\end{adjustbox}
\end{center}

\begin{boxresponda}
\textbf{Responda:} Com base na dinâmica apresentada, identifique qual conjunto de fenômenos naturais (erosão, fusão, solidificação ou compactação) está oculto sob o rótulo \textbf{Processo X} para permitir a formação das rochas sedimentares a partir das ígneas.
\end{boxresponda}
\par\vspace{\espacoQuestao}

\noindent\textcolor{CorLinha}{\rule{\linewidth}{0.5pt}} 
\par\vspace{\espacoPreQuestao}
\subsubsection*{14.} A classificação das rochas baseia-se na sua gênese (origem). Julgue as afirmativas a seguir, classificando-as como Verdadeiras (V) ou Falsas (F).

\begin{enumerate}[label=(\ \ ), itemsep=\espacoAlt]
\item As rochas ígneas ou magmáticas, como o granito e o basalto, formam-se exclusivamente pelo resfriamento e solidificação do magma.
\item O mármore é classificado como uma rocha sedimentar, originada pela deposição de areia no fundo de antigos oceanos.
\item Rochas metamórficas surgem quando rochas pré-existentes sofrem transformações estruturais devido a altas pressões e temperaturas, sem chegarem a derreter.
\item A compactação de fragmentos de outras rochas, restos de animais e vegetais dá origem às rochas sedimentares.
\end{enumerate}
\par\vspace{\espacoQuestao}

\noindent\textcolor{CorLinha}{\rule{\linewidth}{0.5pt}} 
\par\vspace{\espacoPreQuestao}
\subsubsection*{15.} Os paleontólogos dedicam suas vidas a escavar e analisar fósseis, mas eles não procuram essas peças geológicas em qualquer tipo de terreno. Justifique, com base na forma de origem, por que é extremamente comum encontrar fósseis em rochas sedimentares, mas praticamente impossível encontrá-los em rochas magmáticas.
\par\vspace{\espacoQuestao}

\noindent\textcolor{CorLinha}{\rule{\linewidth}{0.5pt}} 
\par\vspace{\espacoPreQuestao}
\subsubsection*{16.} Organize as ideias: Associe os exemplos clássicos de rochas aos seus respectivos grupos de classificação estrutural.

\begin{enumerate}[label=\Roman*., itemsep=\espacoAlt]
\item Rocha originada do resfriamento rápido da lava na superfície (Ex: Basalto).
\item Rocha formada pela alteração do calcário submetido a enorme pressão no interior da crosta (Ex: Mármore).
\item Rocha estruturada pela aglomeração e cimentação de grãos de quartzo (Ex: Arenito).
\end{enumerate}
\par\vspace{\espacoQuestao}

\noindent\textcolor{CorLinha}{\rule{\linewidth}{0.5pt}} 
\par\vspace{\espacoPreQuestao}
\subsubsection*{17.} O perfil estratigráfico revela uma encosta escavada pela ação de um rio ao longo de milênios, expondo diferentes leitos rochosos.

\begin{center}
% ==========================================
% CONTROLE INDIVIDUAL DESTA IMAGEM:
% ==========================================
\begin{adjustbox}{trim=0cm 0cm 0cm 0cm, clip, max width=\larguraTikz}
\begin{tikzpicture}
\definecolor{camada1}{HTML}{D7CCC8}
\definecolor{camada2}{HTML}{A1887F}
\definecolor{camada3}{HTML}{795548}
\definecolor{camada4}{HTML}{4E342E}
\definecolor{fundo}{HTML}{E3F2FD}

\fill[fundo, rounded corners=5pt, drop shadow={opacity=0.15, shadow xshift=2pt, shadow yshift=-2pt}] (-1,0) rectangle (6,5);

\fill[camada4, drop shadow={opacity=0.3}] (0,0.5) rectangle (4.5,1.3);
\fill[camada3, drop shadow={opacity=0.3}] (0,1.3) rectangle (4.3,2.2);
\fill[camada2, drop shadow={opacity=0.3}] (0,2.2) rectangle (4.1,3.2);
\fill[camada1, drop shadow={opacity=0.3}] (0,3.2) rectangle (3.9,4.2);

\fill[green!50!black, rounded corners=1pt] (0,4.2) rectangle (3.9,4.4);

\node[font=\sffamily\footnotesize, white, align=center] at (2,0.9) {Camada W};
\node[font=\sffamily\footnotesize, white, align=center] at (2,1.75) {Camada X};
\node[font=\sffamily\footnotesize, white, align=center] at (2,2.7) {Camada Y};
\node[font=\sffamily\footnotesize, black, align=center] at (2,3.7) {Camada Z};

\draw[thick, -stealth] (-0.5, 4.4) -- (-0.5, 0.5);
\node[font=\sffamily\tiny, rotate=90, align=center] at (-0.8, 2.5) {Aprofundamento Escavação};
\end{tikzpicture}
\end{adjustbox}
\end{center}

\begin{boxresponda}
\textbf{Responda:} Aplicando o Princípio da Superposição das camadas geológicas (considerando que não houve dobramentos ou inversões tectônicas nesse terreno), determine qual das camadas (W, X, Y ou Z) é a mais antiga e justifique fisicamente a sua resposta.
\end{boxresponda}
\par\vspace{\espacoQuestao}

\noindent\textcolor{CorLinha}{\rule{\linewidth}{0.5pt}} 
\par\vspace{\espacoPreQuestao}
\subsubsection*{18.} (Estilo VUNESP) O Pão de Açúcar, um dos mais famosos cartões-postais da cidade do Rio de Janeiro, é formado predominantemente por gnaisse, uma rocha que apresenta bandas claras e escuras intercaladas. Essa rocha não se originou de erupções vulcânicas diretas nem da deposição em leitos de rios, mas sim de granitos antiquíssimos que foram submetidos a condições extremas nas profundezas da crosta ao longo de eras geológicas.

Considerando os processos geológicos descritos, a rocha que forma o Pão de Açúcar classifica-se como:

\begin{enumerate}[label=\alph*), itemsep=\espacoAlt]
\item Sedimentar química, resultante da evaporação de oceanos primitivos.
\item Magmática intrusiva, resfriada lentamente no núcleo da Terra.
\item Metamórfica, resultante da alteração estrutural promovida por calor e pressão.
\item Magmática efusiva, originada pelo rápido resfriamento da lava no mar.
\item Sedimentar clástica, formada pela compactação da areia da praia ao longo dos milênios.
\end{enumerate}
\par\vspace{\espacoQuestao}

\noindent\textcolor{CorLinha}{\rule{\linewidth}{0.5pt}} 
\par\vspace{\espacoPreQuestao}
\subsubsection*{19.} (ENEM - Adaptada) Monumentos históricos de pedra calcária ou mármore, localizados em grandes centros urbanos industrializados, sofrem um desgaste muito mais acelerado do que aqueles situados em regiões rurais. A interação dos gases poluentes (como dióxido de enxofre e óxidos de nitrogênio) com a umidade atmosférica gera soluções ácidas que reagem quimicamente com os minerais dessas esculturas, desgastando-as progressivamente.

Esse processo de alteração e desgaste das rochas pela interação com agentes atmosféricos e químicos é denominado:

\begin{enumerate}[label=\alph*), itemsep=\espacoAlt]
\item Vulcanismo primário.
\item Intemperismo.
\item Tectonismo de placas.
\item Litificação.
\item Cristalização magmática.
\end{enumerate}
\par\vspace{\espacoQuestao}

\subsection*{Capítulo 4: Fósseis}
\vspace{-0.2cm}\noindent\rule{\linewidth}{1.5pt} 
\par\vspace{\espacoPosTitulo}
\subsubsection*{20.} O diagrama esquematiza o processo tafonômico (fossilização) de um organismo marinho, desde o instante da morte até a sua descoberta geológica.

\begin{center}
% ==========================================
% CONTROLE INDIVIDUAL DESTA IMAGEM:
% ==========================================
\begin{adjustbox}{trim=0cm 0cm 0cm 0cm, clip, max width=\larguraTikz}
\begin{tikzpicture}
\definecolor{fundo}{HTML}{FDFEFE}
\definecolor{agua}{HTML}{AED6F1}
\definecolor{areia}{HTML}{F5CBA7}
\definecolor{rocha1}{HTML}{D4AC0D}
\definecolor{rocha2}{HTML}{9A7D0A}

\fill[fundo, rounded corners=5pt, drop shadow={opacity=0.15, shadow xshift=2pt, shadow yshift=-2pt}] (-1,0) rectangle (13,4);

\fill[agua, opacity=0.7] (0,1) rectangle (3.5,3.5);
\fill[areia] (0,0.5) rectangle (3.5,1);
\node[font=\sffamily\scriptsize, align=center] at (1.75, 3.7) {1. Morte do Organismo};
\draw[thick, gray] (1.75, 2.5) ellipse (0.4 and 0.15);
\draw[thick, gray] (1.35, 2.5) -- (1.1, 2.7) -- (1.1, 2.3) -- cycle;
\draw[-stealth, dashed, gray] (1.75, 2.2) -- (1.75, 1.2);

\fill[agua, opacity=0.7] (4.5,2.5) rectangle (8,3.5);
\fill[rocha1] (4.5,1.5) rectangle (8,2.5);
\fill[areia] (4.5,0.5) rectangle (8,1.5);
\node[font=\sffamily\scriptsize, align=center] at (6.25, 3.7) {2. Deposição de Sedimentos};
\draw[thick, black] (6.25, 1) -- (6.6, 1); \draw[thick, black] (6.4, 0.9) -- (6.4, 1.1);
\node[font=\sffamily\tiny, align=center] at (6.25, 2) {Acúmulo \\ de Areia/Lama};

\fill[rocha2] (9,2) rectangle (12.5,3.5);
\fill[rocha1] (9,1.5) rectangle (12.5,2);
\fill[areia] (9,0.5) rectangle (12.5,1.5);
\node[font=\sffamily\scriptsize, align=center] at (10.75, 3.7) {3. Compactação (Milhões de anos)};
\draw[ultra thick, white] (10.75, 1) -- (11.1, 1); \draw[ultra thick, white] (10.9, 0.9) -- (10.9, 1.1);

\draw[thick, dashed, gray] (4,0.5) -- (4,3.5);
\draw[thick, dashed, gray] (8.5,0.5) -- (8.5,3.5);
\end{tikzpicture}
\end{adjustbox}
\end{center}

\begin{boxresponda}
\textbf{Responda:} Com base na sequência ilustrada, explique qual é a principal razão pela qual o "soterramento rápido" (Painel 2) por lama ou areia é uma condição crucial para que os restos do animal se transformem em fóssil, em vez de desaparecerem completamente.
\end{boxresponda}
\par\vspace{\espacoQuestao}

\noindent\textcolor{CorLinha}{\rule{\linewidth}{0.5pt}} 
\par\vspace{\espacoPreQuestao}
\subsubsection*{21.} A formação de um fóssil é considerada um "acidente geológico", pois a maioria absoluta dos seres vivos que já existiram na Terra não deixou vestígios. Julgue as afirmativas como Verdadeiras (V) ou Falsas (F).

\begin{enumerate}[label=(\ \ ), itemsep=\espacoAlt]
\item Partes moles dos organismos, como músculos, pele e órgãos internos, são facilmente fossilizadas na maioria das rochas sedimentares.
\item Pegadas, marcas de cauda e fezes petrificadas (coprólitos) são classificadas como icnofósseis, ou fósseis de vestígios.
\item O âmbar, uma resina vegetal fossilizada, tem a capacidade de preservar insetos inteiros, mantendo intactos detalhes de seus corpos.
\item Para datar as eras geológicas da Terra, os paleontólogos preferem utilizar rochas magmáticas formadas na atualidade, pois contêm os fósseis mais preservados.
\end{enumerate}
\par\vspace{\espacoQuestao}

\noindent\textcolor{CorLinha}{\rule{\linewidth}{0.5pt}} 
\par\vspace{\espacoPreQuestao}
\subsubsection*{22.} Os fósseis não são apenas curiosidades de museu; eles funcionam como cápsulas do tempo da biologia e da geologia. Descreva, com suas palavras, qual a importância do estudo dos fósseis para a compreensão das eras geológicas e da evolução da vida em nosso planeta.
\par\vspace{\espacoQuestao}

\noindent\textcolor{CorLinha}{\rule{\linewidth}{0.5pt}} 
\par\vspace{\espacoPreQuestao}
\subsubsection*{23.} Organize as ideias: Associe os tipos de registro paleontológico com a sua correta definição.

\begin{enumerate}[label=\Roman*., itemsep=\espacoAlt]
\item Fóssil de Corpo (Somatofóssil).
\item Fóssil de Vestígio (Icnofóssil).
\item Molde Fóssil.
\end{enumerate}

\begin{boxresponda}
\textbf{Associe corretamente:} \\
(\ \ ) O organismo original foi totalmente dissolvido, mas deixou uma cavidade na rocha com o formato exato do seu corpo. \\
(\ \ ) Preservação de ossos, dentes, carapaças ou conchas de um organismo. \\
(\ \ ) Marcas de atividade vital preservadas em rocha, como pegadas de um dinossauro na lama que petrificou.
\end{boxresponda}
\par\vspace{\espacoQuestao}

\noindent\textcolor{CorLinha}{\rule{\linewidth}{0.5pt}} 
\par\vspace{\espacoPreQuestao}
\subsubsection*{24.} (Estilo OBA) Na bioestratigrafia, fósseis específicos podem ser usados para correlacionar a idade de rochas em diferentes partes do mundo. O esquema ilustra perfis rochosos de duas montanhas situadas a centenas de quilômetros de distância.

\begin{center}
% ==========================================
% CONTROLE INDIVIDUAL DESTA IMAGEM:
% ==========================================
\begin{adjustbox}{trim=0cm 0cm 0cm 0cm, clip, max width=\larguraTikz}
\begin{tikzpicture}
\definecolor{fundo}{HTML}{FAFAFA}
\definecolor{camA}{HTML}{D2B4DE}
\definecolor{camB}{HTML}{A9CCE3}
\definecolor{camC}{HTML}{F9E79F}
\definecolor{camD}{HTML}{F5CBA7}

\fill[fundo, rounded corners=5pt, drop shadow={opacity=0.15, shadow xshift=2pt, shadow yshift=-2pt}] (-1,0) rectangle (9,5.5);

\node[font=\sffamily\bfseries, align=center] at (1.5, 5) {Sítio Arqueológico 1};
\fill[camD, drop shadow={opacity=0.2}] (0,0.5) rectangle (3,1.5);
\fill[camC, drop shadow={opacity=0.2}] (0,1.5) rectangle (3,2.5);
\fill[camB, drop shadow={opacity=0.2}] (0,2.5) rectangle (3,3.5);
\fill[camA, drop shadow={opacity=0.2}] (0,3.5) rectangle (3,4.5);

\node[font=\sffamily\bfseries, align=center] at (6.5, 5) {Sítio Arqueológico 2};
\fill[camC, drop shadow={opacity=0.2}] (5,0.5) rectangle (8,1.5);
\fill[camB, drop shadow={opacity=0.2}] (5,1.5) rectangle (8,2.5);
\fill[camA, drop shadow={opacity=0.2}] (5,2.5) rectangle (8,3.5);
\fill[gray!30, drop shadow={opacity=0.2}] (5,3.5) rectangle (8,4.5);

\node[font=\Huge, align=center] at (1.5, 2) {\textbf{$\sim$}};
\node[font=\sffamily\tiny, align=center] at (1.5, 1.8) {Fóssil Guia};

\node[font=\Huge, align=center] at (6.5, 1) {\textbf{$\sim$}};
\node[font=\sffamily\tiny, align=center] at (6.5, 0.8) {Fóssil Guia};

\node[fill=white, inner sep=1pt, rounded corners=1pt, font=\sffamily\tiny, align=center] at (0.5, 4) {Cam. 1};
\node[fill=white, inner sep=1pt, rounded corners=1pt, font=\sffamily\tiny, align=center] at (0.5, 3) {Cam. 2};
\node[fill=white, inner sep=1pt, rounded corners=1pt, font=\sffamily\tiny, align=center] at (0.5, 2) {Cam. 3};
\node[fill=white, inner sep=1pt, rounded corners=1pt, font=\sffamily\tiny, align=center] at (0.5, 1) {Cam. 4};

\node[fill=white, inner sep=1pt, rounded corners=1pt, font=\sffamily\tiny, align=center] at (5.5, 4) {Cam. W};
\node[fill=white, inner sep=1pt, rounded corners=1pt, font=\sffamily\tiny, align=center] at (5.5, 3) {Cam. X};
\node[fill=white, inner sep=1pt, rounded corners=1pt, font=\sffamily\tiny, align=center] at (5.5, 2) {Cam. Y};
\node[fill=white, inner sep=1pt, rounded corners=1pt, font=\sffamily\tiny, align=center] at (5.5, 1) {Cam. Z};
\end{tikzpicture}
\end{adjustbox}
\end{center}

\begin{boxresponda}
\textbf{Responda:} Observando a presença do mesmo "Fóssil Guia" em ambos os locais, qual camada do Sítio Arqueológico 2 (W, X, Y ou Z) possui aproximadamente a mesma idade geológica da \textbf{Camada 3} do Sítio Arqueológico 1? 
\end{boxresponda}
\par\vspace{\espacoQuestao}

\noindent\textcolor{CorLinha}{\rule{\linewidth}{0.5pt}} 
\par\vspace{\espacoPreQuestao}
\subsubsection*{25.} (Estilo ENEM) Para que um fóssil seja considerado um excelente "fóssil-guia" ou "fóssil-índice" e auxilie os cientistas na datação de estratos rochosos em escala global, a espécie extinta deve possuir algumas características vitais. Um exemplo notável são os Trilobitas, artrópodes marinhos que dominaram os mares da Era Paleozoica.

Para ser um fóssil-guia eficiente, a espécie fossilizada deve ter apresentado em vida:

\begin{enumerate}[label=\alph*), itemsep=\espacoAlt]
\item Existência prolongada por todas as eras geológicas e distribuição territorial muito restrita.
\item Existência num período de tempo geológico relativamente curto e ampla distribuição geográfica mundial.
\item Corpo predominantemente formado por tecidos moles e ausência total de carapaças ou esqueletos.
\item Adaptação exclusiva a ambientes vulcânicos de altíssimas temperaturas.
\item Rápida decomposição imediata após a morte, impedindo a mineralização de seus ossos.
\end{enumerate}
\par\vspace{\espacoQuestao}

\end{multicols*}
\end{document}